\documentclass{beamer}
\usetheme{Berlin}

\title{Systems with a Small Number of
Variables}
\subtitle{Course: Complex Systems Modeling}
\author{Robert Flassig}
\institute{Brandenburg University of Applied Sciences}
\date{\today}

\begin{document}
\begin{frame}
\titlepage
\end{frame}

\begin{frame}
\frametitle{Outline}
\tableofcontents
\end{frame}

\section{Continuous-Time Dynamical System}

\begin{frame}
\frametitle{Continuous-Time Models with Differential Equations 1/2}

\begin{itemize}
	\scriptsize
    \item Continuous-time models are described using differential equations.
    \item These models are prevalent in science and engineering due to their ability to represent natural phenomena occurring smoothly over time (e.g., motion of objects, flow of electric current).
\end{itemize}

\begin{block}{General Formulation of a First-Order Continuous-Time Model}
\[
\frac{dx}{dt} = F(x, t)
\]
\begin{itemize}
    \item \(x\): State of the system (scalar or vector variable).
    \item \(\frac{dx}{dt}\): Time derivative of \(x\), formally defined as:
    \[
    \frac{dx}{dt} = \lim_{\delta t \to 0} \frac{x(t + \delta t) - x(t)}{\delta t}
    \]
\end{itemize}
\end{block}

\end{frame}

\begin{frame}
\frametitle{Continuous-Time Models with Differential Equations 2/2}

\begin{itemize}
    \item Integrating the model over time \(t\) provides the system's trajectory.
    \item Integration can sometimes be performed algebraically, but computational simulation (numerical integration) is often used in practice.
    \item A fundamental assumption:
    \begin{itemize}
        \item Trajectories of the system's state are smooth everywhere in the phase space.
        \item Instantaneous abrupt changes, possible in discrete-time models, do not occur in continuous-time models.
    \end{itemize}
    $\rightarrow$ Recall discussion on the assumptions of the  predator-prey model, e.g. $1.4$ rabbit, $2.5$ wolves. Also small vs. large number of entities - continuum assumption.
\end{itemize}

\end{frame}

% Frame 1: Introduction to Classifications of Model Equations
\begin{frame}
\frametitle{Classifications of Model Equations}

\begin{itemize}
    \item We may distinct between \underline{linear} and \underline{nonlinear systems}, as well as \underline{autonomous} and \underline{non-autonomous systems}.
    \begin{itemize}
        \item \underline{First-order systems}: Involve only first-order derivatives (\(\frac{dx}{dt}\)).
        \item \underline{Higher-order systems}: Involve higher-order derivatives (\(\frac{d^2x}{dt^2}, \frac{d^3x}{dt^3}, \ldots\)).
    \end{itemize}
    \item Non-autonomous, higher-order systems can always be transformed into:
    \begin{itemize}
        \item Autonomous, first-order forms.
        \item This is achieved by introducing additional state variables.
    \end{itemize}
\end{itemize}

\end{frame}



\begin{frame}
\frametitle{Classification of Systems: Linear vs. Nonlinear}
\begin{columns}[T,onlytextwidth]
    \column{0.48\textwidth}
    \textbf{Linear Systems}
    \begin{itemize}
        \item \(\dot{\mathbf{x}} = A\mathbf{x} + B\mathbf{u}\)
        \item Scalar: \(\dot{x} = a x + b u\)
        \item Superposition principle holds
    \end{itemize}

    \column{0.48\textwidth}
    \textbf{Nonlinear Systems}
    \begin{itemize}
        \item \(\dot{\mathbf{x}} = f(\mathbf{x}, \mathbf{u})\)
        \item Scalar: \(\dot{x} = a x^2 + b u\)
        \item Complex dynamics: chaos, bifurcations, limit cycles
    \end{itemize}
\end{columns}
\end{frame}




\begin{frame}
\frametitle{Classification of Systems: Autonomous vs. Non-Autonomous}
\scriptsize
\begin{columns}[T,onlytextwidth]
    %-------------------------------------
    \column{0.48\textwidth}
    \textbf{Autonomous Systems}
    \begin{itemize}
        \item Evolution depends only on the state, not explicitly on time.
        \item \textbf{Matrix form:}
        \[
        \dot{\mathbf{x}} = A\mathbf{x}
        \]
        \item \textbf{Scalar example:}
        \[
        \frac{dx}{dt} = a x
        \]
        \item System is \textit{time-invariant} → suitable for phase-space or equilibrium analysis.
    \end{itemize}

    %-------------------------------------
    \column{0.48\textwidth}
    \textbf{Non-Autonomous Systems}
    \begin{itemize}
        \item Explicit time-dependence in the equations or inputs.
        \item \textbf{Matrix form:}
        \[
        \dot{\mathbf{x}} = A(t)\mathbf{x} + B(t)\mathbf{u}(t)
        \]
        \item \textbf{Scalar example:}
        \[
        \frac{dx}{dt} = a(t) x + b(t)
        \]
        \item Require special treatment for time-dependent forcing or varying parameters.
    \end{itemize}
\end{columns}
\end{frame}



\begin{frame}
\frametitle{Example: Pendulum Dynamics}
\scriptsize
\begin{columns}[T,onlytextwidth]

  %---------------------------------------
  \column{0.56\textwidth}
  \textbf{Second-order (nonlinear) equation of motion:}
  \[
    \frac{d^2\theta}{dt^2} = -\frac{g}{L}\sin\theta
  \]

  \textbf{Interpretation:}
  \begin{itemize}\itemsep2pt
    \item The restoring torque is proportional to \(-\sin\theta\).
    \item Equation is second-order and nonlinear.
    \item For large angles, motion deviates from simple harmonic behavior.
  \end{itemize}

  %---------------------------------------
  \column{0.44\textwidth}
  \textbf{Variables}
  \vspace{-4pt}
  \begin{itemize}\itemsep2pt
    \item \( \theta \): angular position
    \item \( g \): gravitational acceleration
    \item \( L \): pendulum length
  \end{itemize}

  \vspace{2pt}
  \textbf{Remarks}
  \vspace{-4pt}
  \begin{itemize}\itemsep2pt
    \item Nonlinear due to \( \sin\theta \)
    \item Linearization possible for small angles
  \end{itemize}

\end{columns}
\end{frame}

\begin{frame}
\frametitle{Converting to a First-Order System}
\scriptsize
\begin{columns}[T,onlytextwidth]

  %---------------------------------------
  \column{0.56\textwidth}
  Introduce angular velocity:
  \[
    \omega = \frac{d\theta}{dt}
  \]

  Substitute into the second-order equation:
  \[
    \begin{aligned}
      \dot{\theta} &= \omega, \\
      \dot{\omega} &= -\frac{g}{L}\sin\theta
    \end{aligned}
  \]

  Compactly written as:
  \[
    \dot{\mathbf{x}} =
    \begin{bmatrix}
      \omega \\[2pt]
      -\dfrac{g}{L}\sin\theta
    \end{bmatrix},
    \quad
    \mathbf{x} =
    \begin{bmatrix}
      \theta \\[2pt]
      \omega
    \end{bmatrix}
  \]

  %---------------------------------------
  \column{0.44\textwidth}
  \textbf{Notes}
  \vspace{-4pt}
  \begin{itemize}\itemsep2pt
    \item Two coupled first-order ODEs
    \item Nonlinear due to \( \sin\theta \)
    \item Enables phase-plane analysis and simulation
  \end{itemize}

\end{columns}
\end{frame}
\begin{frame}
\frametitle{Linearized Pendulum (Small Angles)}
\scriptsize

\begin{columns}[T,onlytextwidth]

  %---------------------------------------
  \column{0.56\textwidth}
  For small angles, approximate \(\sin\theta \approx \theta\):
  \[
    \begin{aligned}
      \dot{\theta} &= \omega,\\[-2pt]
      \dot{\omega} &= -\frac{g}{L}\theta
    \end{aligned}
  \]

  \textbf{Matrix form:}
  \[
    \dot{\mathbf{x}} =
    \underbrace{\begin{bmatrix}
      0 & 1\\[2pt]
      -\dfrac{g}{L} & 0
    \end{bmatrix}}_{A}
    \mathbf{x}
  \]

  %---------------------------------------
  \column{0.44\textwidth}
  \textbf{Remarks}
  \vspace{-4pt}
  \begin{itemize}\itemsep2pt
    \item Linear time-invariant (LTI) system
    \item Represents a \textit{simple harmonic oscillator}
    \item Solution:
      \[
        \theta(t) = \theta_0 \cos(\sqrt{g/L}\,t)
      \]
    \item Basis for control and small-signal stability analysis
  \end{itemize}

\end{columns}
\end{frame}



\begin{frame}
\frametitle{Conversion of Higher-Order Equations}

\begin{itemize}
    \item The technique for converting higher-order equations to first-order forms works for:
    \begin{itemize}
        \item 3rd/4th/.../-order equations,
        \item Any higher-order equations (as long as the highest order is finite).
    \end{itemize}
    \item This method ensures equations are represented as autonomous, first-order systems for analysis and simulation.
\end{itemize}

\begin{block}{Key Advantage}
Autonomous, first-order forms are universal and cover all dynamics of higher-order and non-autonomous systems.
\end{block}

\end{frame}

\section{Example}

\begin{frame}

\frametitle{A Non-Autonomous Driven Pendulum}

\begin{block}{Second-Order Non-Autonomous Equation}
\[
\frac{d^2\theta}{dt^2} = -\frac{g}{L} \sin\theta + k \sin(2\pi f t + \phi) \tag{6.7}
\]
\begin{itemize}
    \item Describes a driven pendulum with:
    \begin{itemize}
        \item \(\frac{g}{L}\): Gravitational term.
        \item \(k \sin(2\pi f t + \phi)\): Periodic external force, e.g., from an electromagnet.
    \end{itemize}
    \item This equation is nonlinear and non-autonomous due to explicit time-dependence (\(t\)).
\end{itemize}
\end{block}

\begin{itemize}
    \item Goal: Convert it into a first-order, autonomous system.
\end{itemize}

\end{frame}

\begin{frame}
\frametitle{Step 1: First-Order Conversion}

\begin{block}{Introduce \(\omega = \frac{d\theta}{dt}\)}
\begin{align}
\frac{d\theta}{dt} &= \omega \tag{6.8} \\
\frac{d\omega}{dt} &= -\frac{g}{L} \sin\theta + k \sin(2\pi f t + \phi) \tag{6.9}
\end{align}
\end{block}

\begin{itemize}
    \item \(\omega\): Angular velocity of the pendulum.
    \item The system is now in first-order form but still non-autonomous due to explicit \(t\) in the sine term.
\end{itemize}

\end{frame}

\begin{frame}
\frametitle{Step 2: Eliminating Time Dependency}
\scriptsize
\begin{columns}[T,onlytextwidth]

  %---------------------------------------
  \column{0.56\textwidth}
  \textbf{Introduce a clock variable \(\tau\):}
  \[
    \frac{d\tau}{dt} = 1, \qquad \tau(0)=0
  \]
  \vspace{-4pt}
  \begin{itemize}\itemsep1.5pt
    \item Ensures \(\tau(t) = t\).
    \item Replaces explicit time dependence with state variable \(\tau\).
    \item Converts a non-autonomous system into an autonomous one.
  \end{itemize}

  \vspace{4pt}
  \textbf{Original (time-dependent) term:}
  \[
    k\sin(2\pi f t + \phi)
  \]
  becomes
  \[
    k\sin(2\pi f \tau + \phi)
  \]

  %---------------------------------------
  \column{0.44\textwidth}
  \textbf{Resulting autonomous system:}
  \[
  \begin{aligned}
    \dot{\theta} &= \omega,\\[-2pt]
    \dot{\omega} &= -\frac{g}{L}\sin\theta 
                   + k\sin(2\pi f\tau + \phi),\\[-2pt]
    \dot{\tau}   &= 1
  \end{aligned}
  \]
  \vspace{-2pt}
  \textbf{Notes}
  \vspace{-3pt}
  \begin{itemize}\itemsep1.5pt
    \item System now autonomous in \((\theta,\omega,\tau)\).
    \item Useful for phase-space analysis and numerical integration.
  \end{itemize}

\end{columns}
\end{frame}

\begin{frame}
\frametitle{Key Insights from the Conversion Process}

\begin{itemize}
    \item Any higher-order, non-autonomous system can be converted into a first-order, autonomous system by:
    \begin{itemize}
        \item Introducing new state variables for higher-order derivatives.
        \item Replacing explicit time dependencies with a clock variable (\(\tau\)).
    \end{itemize}
    \item The resulting equations:
    \begin{itemize}
        \item Are always first-order.
        \item Have no explicit time dependency (autonomous).
    \end{itemize}
\end{itemize}

\begin{block}{General Utility}
This method ensures compatibility with standard tools for analyzing first-order, autonomous systems.
\end{block}

\end{frame}

\section{Exercise}

%------------------------------------------
\begin{frame}
\frametitle{Exercise 1: Convert to First-Order Form}
\scriptsize
\begin{columns}[T,onlytextwidth]

  %---------------------------------------
  \column{0.55\textwidth}
  \textbf{Given Equation:}
  \[
    \frac{d^2x}{dt^2} - x\frac{dx}{dt} + x^2 = 0
  \]

  \textbf{Introduce a new variable:}
  \[
    v = \frac{dx}{dt}
  \]

  Substitute and rewrite:
  \[
    \begin{aligned}
      \dot{x} &= v,\\[-2pt]
      \dot{v} &= x v - x^2
    \end{aligned}
  \]

  %---------------------------------------
  \column{0.45\textwidth}
  \textbf{First-Order System:}
  \[
    \begin{aligned}
      \dot{x} &= v,\\[-2pt]
      \dot{v} &= x v - x^2
    \end{aligned}
  \]

  \vspace{-2pt}
  \textbf{Remarks:}
  \vspace{-4pt}
  \begin{itemize}\itemsep1.5pt
    \item Nonlinear due to product term \(xv\).
    \item Expressed as a system of two coupled ODEs.
    \item Easier to simulate or visualize in the phase plane.
  \end{itemize}

\end{columns}
\end{frame}
%------------------------------------------
\begin{frame}
\frametitle{Exercise 2: Convert to Autonomous, First-Order Form}
\scriptsize
\begin{columns}[T,onlytextwidth]

  %---------------------------------------
  \column{0.55\textwidth}
  \textbf{Given Equation:}
  \[
    \frac{d^2x}{dt^2} - a\cos(bt) = 0
  \]

  \textbf{Step 1: Introduce} \( v = \frac{dx}{dt} \)
  \[
    \dot{x} = v, \qquad \dot{v} = a\cos(bt)
  \]

  \textbf{Step 2: Eliminate explicit time dependency:}
  \[
    \frac{d\tau}{dt} = 1, \quad \tau(0)=0
  \]

  Replace \(t \to \tau\) in the forcing term:
  \[
    \cos(bt) \to \cos(b\tau)
  \]

  %---------------------------------------
  \column{0.45\textwidth}
  \textbf{Autonomous, First-Order Form:}
  \[
    \begin{aligned}
      \dot{x} &= v,\\[-2pt]
      \dot{v} &= a\cos(b\tau),\\[-2pt]
      \dot{\tau} &= 1
    \end{aligned}
  \]

  \vspace{-2pt}
  \textbf{Remarks:}
  \vspace{-4pt}
  \begin{itemize}\itemsep1.5pt
    \item System is now autonomous in \((x,v,\tau)\).
    \item Clock variable replaces explicit time.
    \item Technique generalizes to many driven systems.
  \end{itemize}

\end{columns}
\end{frame}

\section{Summary}

\begin{frame}
\frametitle{Summary: Differential Equations and System Behavior}
\scriptsize
\begin{columns}[T,onlytextwidth]

  %---------------------------------------
  \column{0.55\textwidth}
  \textbf{1. Typical Behaviors of Linear Systems}
  \begin{itemize}\itemsep2pt
    \item Exponential growth or decay
    \item Periodic oscillations
    \item Stationary (equilibrium) states
    \item Combined forms, e.g. exponentially damped oscillations
  \end{itemize}

  \vspace{2pt}
  \textbf{2. Extended Behaviors (Non-Diagonalizable Systems)}
  \begin{itemize}\itemsep2pt
    \item Solutions may include
      \begin{itemize}\itemsep1pt
        \item Polynomials in time
        \item Products of polynomials and exponentials
      \end{itemize}
    \item Occurs when the coefficient matrix cannot be fully diagonalized
  \end{itemize}

  %---------------------------------------
  \column{0.45\textwidth}
  \textbf{3. Solvability and Methods}
  \begin{itemize}\itemsep2pt
    \item Linear ODEs: typically analytically solvable
    \item Nonlinear ODEs: often require numerical integration
    \item Many nonlinear systems exhibit complex phenomena:
      \begin{itemize}\itemsep1pt
        \item Chaos, bifurcations, limit cycles
      \end{itemize}
  \end{itemize}

  \vspace{4pt}
  \begin{block}{Key Insight}
  Understanding these distinctions guides model selection  
  and the choice of analytical vs.\ numerical solution methods.
  \end{block}

\end{columns}
\end{frame}


\end{document}